\section{Further Work: Comparing Implementations of a Primitive}
\label{sec:further}

The deconstruction shows that the five primitives cover what exists: every component in nineteen harnesses had a home in one primitive or one concern. It does not show which implementation of a primitive is better, for which model, on which task. That is the comparison the vocabulary makes well-posed, and it has not been done.

\paragraph{The form of the study.} Hold the model fixed. Hold four primitives fixed. Vary the fifth, one design decision at a time, and measure on a task set with a check that does not depend on the primitive varied. The harness is stated in the five parts of Table~\ref{tab:primitives}, so that the varied part is the only difference between arms, and the statement is published with the result. \citet{fan2026empirical} is the nearest existing study, with three decisions varied across four models; the vocabulary names the remaining decisions and says what must be held fixed.

\paragraph{One comparison per primitive.} The decisions in Table~\ref{tab:decisions} are the candidates. Within \prim{Context}: window fitting, where the harnesses in \S\ref{sec:deconstruction} drop old turns, summarize, compact at a threshold with a verbatim tail, or leave it to the provider, and no two agree on when; and the placement of project instructions, preloaded or searched. Within \prim{Output}: the action format, structured calls against model-written code run in an interpreter, which some of the nineteen offer and the rest do not; and the execution environment, host, OS sandbox or container, which changes what the model is allowed to do and what it is told it did. Within \prim{Control flow}: the stopping rule, a step budget against a cost budget against none; and the approval gate, a rule pipeline against a classifier model against per-tool assertion, measured on both task completion and the rate of harmful actions. Within \prim{Model interface}: the retry and backup policy, measured under induced provider failure. Within \prim{Input}: the delivery of mid-turn messages, at the next step boundary or after the turn, measured on tasks where the person changes the request part-way.

\paragraph{What such a study needs.} A task set with a harness-independent check. A model held at a fixed version, since the harness that suits one model may not suit another \citep{metr2025update}. The four fixed primitives reported in the same five-part schema, so that a later study can hold the same things fixed. And the evidence files of Appendix~\ref{app:evidence} are a starting point for the arms: each cell of Table~\ref{tab:summary} is an implementation that exists and can be ported into a common harness, because the primitives give the seam along which to port it.

\paragraph{These comparisons are also the test of the framework.} The deconstruction tests coverage: can every component be placed? A comparison tests something the deconstruction cannot: can a result be attributed to the primitive varied? If an implementation cannot be placed in one primitive, or if varying one primitive changes behaviour only through another, the framework is wrong in a way that reading source could not show. The five questions of \S\ref{sec:debugging} make the same prediction in the small: a failure isolated to one primitive should be fixed by a change to that primitive alone. Both predictions are checkable, and neither has been checked at scale.

\paragraph{A second coding of the deconstruction.} The assignments in Appendix~\ref{app:evidence} were made by the author with an AI coding agent as a research instrument, against the definitions of \S\ref{sec:primitives} and the rules of Appendix~\ref{app:rules}. No independent second coding was made, so the agreement between two people applying the definitions is unmeasured. The brief, the definitions and the component lists are published for that purpose.
